\documentclass[10pt,a4paper,sans]{moderncv}
\moderncvstyle{classic}
\moderncvcolor{blue}
\usepackage[utf8]{inputenc}
\usepackage[scale=0.88]{geometry}
\usepackage{helvet}

\name{Alexandro}{Disla}
\title{Data Analyste | Développeur d'Applications Web}
\address{14, rue Saint-Louis Jeanty, Route de Frère}{Port-au-Prince, Haïti}{}
\phone[mobile]{(+509) 4148-3700}
\phone[fixed]{(+509) 4630-7530}
\email{alexandrodisla@hotmail.com}
\extrainfo{Né le 10 juillet 1991 à Port-au-Prince --- Nationalité haïtienne --- GitHub : \url{https://github.com/AD0791}}

\begin{document}
\makecvtitle

\section{Profil Professionnel}
\cvitem{}{Professionnel des systèmes d'information de gestion sociale et sanitaire, avec plus de dix ans de pratique répartie entre l'\textbf{administration publique haïtienne} --- près de huit ans au Ministère de la Planification et de la Coopération Externe, en suivi d'exécution du Plan Triennal d'Investissement puis en modélisation de prévision et d'analyse d'impact des politiques publiques --- et les \textbf{programmes financés par des bailleurs internationaux}.\newline\newline Mon profil réunit les deux métiers que recherche l'UGP-MAST, parce que je les exerce sur les mêmes systèmes. Je \textbf{conçois et développe} des applications web de gestion de données (PHP/Laravel, MySQL, PostgreSQL) et je \textbf{traite et analyse} les données qu'elles hébergent : programmation de formulaires XLSForm pour ODK, réception automatisée des soumissions, apurement et dédoublonnage des bases, contrôles de cohérence et de géolocalisation, puis production d'indicateurs et de tableaux de bord pour la décision. J'ai construit et je maintiens en production le \textbf{back-office de suivi-évaluation d'Anseye Pou Ayiti}, une application PHP/MySQL alimentée par des enquêtes ODK --- l'architecture même d'un système d'information de ciblage et de suivi des bénéficiaires.}

\section{Compétences Clés}
\cvitem{}{
\begin{itemize}
    \item \textbf{Développement d'applications web de gestion de données :} Applications PHP en production (Laravel 8 et PHP natif), API REST, panneaux d'administration, gestion des sessions et des droits d'accès, export Excel, sauvegardes automatisées de base de données, déploiement conteneurisé et documentation du code sous contrôle de version ; conception de schémas relationnels, migrations, optimisation de requêtes et pipelines d'intégration entre systèmes hétérogènes.
    \item \textbf{Collecte numérique et administration de serveurs d'enquête :} Programmation de questionnaires \textbf{XLSForm pour ODK, KoboCollect, CommCare et mWater} --- logiques de saut, contraintes de validation, listes en cascade, collecte hors ligne ; administration du serveur et du panneau d'administration ODK, réception automatisée des soumissions par webhook et intégration directe en base relationnelle.
    \item \textbf{Qualité, traçabilité et consolidation des données :} Nettoyage, \textbf{dédoublonnage}, harmonisation de structures hétérogènes et fusion de jeux de données historiques ; protocoles de Data Quality Assessment établis et supervisés ; contrôles automatisés de complétude, de cohérence, de métadonnées (horodatage, identifiants, géolocalisation) et audits d'enregistrements orphelins, avec suivi des corrections jusqu'à clôture.
    \item \textbf{Analyse statistique et restitution :} Statistique descriptive et inférentielle, plans d'échantillonnage, pondération et désagrégation, analyse de tendances et d'écarts, prévision et analyse d'impact de politiques publiques ; tableaux de bord et rapports de synthèse destinés à des lecteurs non techniciens.
    \item \textbf{Transfert de compétences :} Formation d'équipes de collecte et de cadres sur les outils numériques, les protocoles de gestion de données et les tableaux de bord ; rédaction de guides opérationnels et de documentation normative.
\end{itemize}}

\section{Systèmes d'Information, Données et Suivi-Évaluation}

\cventry{Jan 2026 -- Mars 2026}{Analyste de Données \& Consultant Suivi-Évaluation}{Anseye Pou Ayiti (APA)}{Haïti}{}{
\begin{itemize}
    \item \textbf{Système d'information de suivi des bénéficiaires :} Développement et maintenance du back-office M\&E de l'organisation --- application web \textbf{PHP/MySQL} hébergée en production, alimentée par trois formulaires \textbf{XLSForm} déployés sur ODK/ONA.io dont les soumissions sont reçues par webhook et intégrées automatiquement en base (31 tables).
    \item \textbf{Identification et dédoublonnage :} Conception d'une clé d'identification unique et déterministe des bénéficiaires (cohorte, initiales, sexe, année de naissance, séquence), avec pages de \textbf{recherche de doublons}, d'audit des enregistrements orphelins et de réconciliation entre la base de recensement et les fiches de suivi périodique.
    \item \textbf{Qualité et continuité de service :} Modules de correction et de suppression contrôlées, export Excel, \textbf{sauvegarde complète de la base au format ZIP} téléchargeable ou expédiée par courriel, journalisation des soumissions, sessions sécurisées et verrouillage automatique par inactivité.
    \item \textbf{Indicateurs et restitution :} Mise en œuvre d'indicateurs de performance pour le suivi des programmes de leadership (AL, PL, DL, EK), tableaux de bord de pilotage et maintenance des scripts d'intégration Python sur AWS.
    \item \textbf{Formation :} Accompagnement des équipes de terrain sur l'usage des outils de collecte numérique et sur les procédures de contrôle qualité.
\end{itemize}}

\cventry{Oct 2023 -- Jan 2026}{Consultant mWater \& Analyste de Données}{HANWASH}{Télétravail}{}{
\begin{itemize}
    \item \textbf{Enquêtes géospatiales :} Déploiement d'enquêtes complexes sur mWater avec relevé et validation de la \textbf{géolocalisation} des points de collecte, et résolution des anomalies remontées du terrain.
    \item \textbf{Dispositif de suivi :} Raffinement du plan de suivi-évaluation et du tableau de suivi des indicateurs (ITT) alignés sur les normes JMP.
    \item \textbf{Restitution et formation :} Tableaux de bord automatisés (mWater, Power BI) connectés aux API de collecte ; guides techniques et renforcement des capacités des équipes.
\end{itemize}}

\cventry{Oct 2021 -- Août 2024}{Officier de Suivi \& Évaluation}{Impact Youth Project, Caris Foundation International}{Haïti}{}{
\begin{itemize}
    \item \textbf{Systèmes de collecte multi-sites :} Mise en place de formulaires numériques (\textbf{CommCare}, HIVHaiti) sur un programme de santé déployé sur plusieurs sites, avec standardisation des formulaires et des flux entre sites.
    \item \textbf{Assurance qualité (DQA) :} Établissement et supervision des protocoles de Data Quality Assessment sur des bases \textbf{MySQL} intégrées, avec remontée des chiffres rapportés aux registres sources et suivi des actions correctives.
    \item \textbf{Reporting vers bailleur :} Production et validation d'indicateurs rapportés à l'USAID via la plateforme nationale DATIM/DHIS2, en qualité de producteur de données.
    \item \textbf{Automatisation et tableaux de bord :} Routines Python et SQL remplaçant des exports-imports manuels, réduisant le temps de traitement de 40 \% ; tableaux de bord (Power BI, Quarto) pour le suivi hebdomadaire des indicateurs.
\end{itemize}}

\cventry{Nov 2015 -- Août 2023}{Économiste-Planificateur --- Direction de Planification Économique et Sociale (DPES)}{Ministère de la Planification et de la Coopération Externe (MPCE)}{Port-au-Prince}{}{
\begin{itemize}
    \item \textbf{Suivi de programmes publics :} Contribution aux rapports de suivi de l'exécution du Plan Triennal d'Investissement 2014--2016 --- écarts entre programmation et réalisation, facteurs de retard.
    \item \textbf{Modélisation et analyse d'impact :} Affecté à l'unité statistique et modélisation, participation à la construction du << Modèle de Planification du Développement >>, outil de prévision, de simulation et d'analyse d'impact des politiques publiques, sous la supervision de Dr. Edouard Nsimba, expert international en modélisation.
    \item \textbf{Cadrage et rédaction :} Réunions de cadrage macroéconomique avec le Ministère de l'Économie et des Finances (MEF), suivi des indicateurs macroéconomiques et rédaction pour des rapports de politiques publiques.
\end{itemize}}

\section{Ingénierie Logicielle et Développement Web}

\cventry{Févr 2026 -- Présent}{Ingénieur Logiciel (Fullstack)}{Tekkod LLC}{Télétravail}{}{
\begin{itemize}
    \item \textbf{Backend PHP/Laravel :} Maintenance et évolution d'un backend \textbf{Laravel} (API utilisateurs, authentification multifacteur, notifications push, panneau d'administration Blade) adossé à \textbf{MySQL}, déployé par conteneur sur Google Cloud Run ; optimisation de requêtes (correction de requêtes N+1) et tests PHPUnit.
    \item \textbf{Intégration inter-systèmes :} Service d'intégration \textbf{Python/FastAPI} sur MongoDB assurant l'échange authentifié avec l'API d'un partenaire externe, avec signature cryptographique des requêtes et documentation OpenAPI.
    \item \textbf{Applications clientes :} Modernisation d'un client mobile React Native (montée de version majeure du framework et de la chaîne de compilation Android/iOS) et développement d'une application web React + TypeScript (Vite).
    \item \textbf{Mise en production :} Les deux applications sont publiées et disponibles sur l'App Store et Google Play ; travail conduit en \textbf{méthode Agile}, avec suivi des tâches, versionnement Git et rapports d'avancement aux parties prenantes.
\end{itemize}}

\cventry{Mars 2021 -- Mai 2024}{Développeur Backend \& Ingénieur de Données}{Tekkod LLC}{Télétravail}{}{
\begin{itemize}
    \item Développement de services web REST sécurisés (FastAPI, JWT, contrôle d'accès par rôles) et application du \textit{repository pattern} ; documentation OpenAPI complète.
    \item Pipelines de traitement de données (ETL, Apache Beam) migrant des données MongoDB vers BigQuery ; administration de bases relationnelles en accès concurrent et optimisation des requêtes.
    \item Infrastructure conteneurisée (Docker), intégration continue et tableaux de bord analytiques.
\end{itemize}}

\cventry{Nov 2015 -- Mars 2021}{Consultant Logiciel \& Spécialiste Intégration de Données}{CassionSoft (Projet UNOPS)}{Haïti}{}{Intégration de données sur un projet des Nations Unies : conception de schémas relationnels, scripts de migration et validation de la logique d'intégration par tests unitaires rigoureux ; systèmes d'authentification et d'autorisation par rôles (OAuth2/JWT) et services REST documentés.}

\section{Formation}
\cvitem{2010--2014}{\textbf{Programme du Diplôme d'Études Supérieures en Économie Appliquée} --- C.T.P.E.A, Centre de Techniques de Planification et d'Économie Appliquée, établissement sous tutelle du MPCE, Port-au-Prince. Cycle régulier de quatre ans en économie et statistique : scolarité complète et ensemble des examens réussis, attestés par le Secrétariat Général du CTPEA. \textit{Attestation jointe au dossier.}}
\cvitem{2009--2010}{Baccalauréat, Philo section C, mention Bien --- Institution Saint-Louis de Gonzague, Port-au-Prince}
\cvitem{Formation continue}{Flatiron School (Software Engineering) ; formations spécialisées en développement web, bases de données et science des données.}

\section{Compétences Techniques}
\cvitem{Langages}{\textbf{PHP} (natif et Laravel), \textbf{SQL}, \textbf{Python}, JavaScript, TypeScript, \textbf{R}}
\cvitem{Web \& Frameworks}{Laravel 8 (Blade, Eloquent), FastAPI, React, React Native, Vite, Tailwind CSS, Bootstrap, API REST, OpenAPI}
\cvitem{Bases de données}{\textbf{MySQL}, \textbf{PostgreSQL}, \textbf{MongoDB}, SQLite, BigQuery --- conception de schémas, migrations, optimisation de requêtes, sauvegarde et restauration}
\cvitem{Serveurs \& Méthodes}{Apache (mod\_rewrite, .htaccess), Tomcat, \textbf{serveur ODK / ONA.io}, Docker, Google Cloud Run, AWS, Azure ; \textbf{Git / GitHub}, \textbf{Agile} (plus de 5 ans), tests automatisés (PHPUnit, pytest), documentation technique et guides utilisateurs}
\cvitem{Collecte de données}{\textbf{ODK / XLSForm}, \textbf{KoboToolbox}, \textbf{CommCare}, ONA.io, mWater, HIVHaiti}
\cvitem{Analyse \& Visualisation}{\textbf{R} (RMarkdown, Quarto), \textbf{Python} (Pandas, NumPy), \textbf{Excel avancé}, \textbf{Power BI}, \textbf{Tableau}, Looker Studio, Chart.js}
\cvitem{Langues}{Créole haïtien (maternel), Français (maternel), Anglais (excellent à l'écrit et à la lecture, professionnel à l'oral)}

\end{document}
